%! Logarithmic Function Context and Data Modeling — Unit 2, Lesson 2.14 — Warm-Up (Answer Key)
\documentclass[10pt]{article}
\usepackage{precalculus-article}
\usepackage{precalculus-key}   % pulls in -boxes; do NOT also load -boxes
\newcommand{\TallMath}[1]{$\displaystyle #1\rule[-1.4em]{0pt}{3.2em}$}

\begin{document}

\pageheader{Unit 2, Lesson 2.14}{Warm-Up --- Answer Key}
\namedateperiod

\vspace{0.15in}

\begin{enumerate}[leftmargin=1.8em, itemsep=16pt]

\item The table below shows a relationship between $x$ and $y$.

\vspace{6pt}
\begin{center}
\renewcommand{\arraystretch}{1.8}
\begin{tabular}{|c|c|c|c|c|}
\hline
\rowcolor{sky}
$x$ & $1$ & $10$ & $100$ & $1000$ \\
\hline
$y$ & $0$ & $1$  & $2$   & $3$ \\
\hline
\end{tabular}
\end{center}

\vspace{4pt}
What type of function does this table suggest? Circle one: \quad
\textbf{Linear} \quad \textbf{Exponential} \quad \ans{Logarithmic}

\vspace{4pt}
Explain how you know: \ansline{As $y$ increases by 1, $x$ multiplies by 10 --- the inputs
change proportionally for equal output increments, the hallmark of a logarithmic model.}

\vspace{4pt}
What is the base of the function? \ans{10}

\item Identify whether each pattern below suggests a \textbf{linear}, \textbf{exponential},
      or \textbf{logarithmic} function. Circle your answer and give a one-sentence reason.

\vspace{6pt}
\begin{enumerate}[leftmargin=1.5em, itemsep=10pt, label=\alph*.]
  \item As $x$ doubles, $y$ increases by a constant amount of 3.

  \vspace{2pt}
  \textbf{Linear} \quad \textbf{Exponential} \quad \ans{Logarithmic} \quad
  because \ansline{equal output increments (add 3) correspond to proportional input changes
  (multiply by 2) --- that is the logarithmic pattern.}

  \item Over equal intervals of $x$, the output $y$ multiplies by 4.

  \vspace{2pt}
  \textbf{Linear} \quad \ans{Exponential} \quad \textbf{Logarithmic} \quad
  because \ansline{outputs change by a constant factor (multiply by 4) over equal input
  intervals --- that is the exponential pattern.}
\end{enumerate}

\item Consider $f(x) = \log_5(x)$.

\begin{enumerate}[leftmargin=1.5em, itemsep=8pt, label=\alph*.]
  \item Evaluate $f(1) = \log_5(1) = $ \ans{0}

  \item The input value where $f(x) = 0$ is called the \textbf{real zero} of $f$.
        What is the real zero of $f(x) = \log_5(x)$? \ans{$x = 1$}

  \item Evaluate $f(5) = \log_5(5) = $ \ans{1}
\end{enumerate}

\end{enumerate}

\end{document}
